% Front matter consolidated here so main.tex only orchestrates inputs.
\pagenumbering{roman}

\chapter*{Table of Contents}
{\singlespacing
\makeatletter\@starttoc{toc}\makeatother
}

\clearpage
\chapter*{Acknowledgements}
\addcontentsline{toc}{chapter}{Acknowledgements}

I would like to express my gratitude to all those who gave me the possibility to complete this thesis, as well as for the academic support I received during my time studying at the University of Science and Technology of Hanoi.

Foremost, I would like to send my sincere thanks to my supervisor, Assoc.\ Prof.\ Nguyen Viet Anh, for his guidance, encouragement, and valuable feedback throughout the course of this internship. His insights helped shape both the direction of the research and the quality of this report.

I would also like to thank my internal supervisor, Dr.\ Giang Anh Tuan, for his continued support and advice during the project.

Finally, I am grateful to the Department of Information and Communication Technology and all the faculty members at USTH whose teaching provided the foundation for the work presented here.

\clearpage
\chapter*{List of Abbreviations}
\addcontentsline{toc}{chapter}{List of Abbreviations}

\begin{longtable}{@{}p{0.22\textwidth}L{0.68\textwidth}@{}}
\toprule
\textbf{Abbreviation} & \textbf{Meaning} \\
\midrule
AI & Artificial Intelligence \\
API & Application Programming Interface \\
BF16 & Brain Floating Point 16-bit \\
CLI & Command-Line Interface \\
CPU & Central Processing Unit \\
CUDA & Compute Unified Device Architecture \\
F1 & Harmonic mean of precision and recall \\
GGUF & GPT-Generated Unified Format \\
GPU & Graphics Processing Unit \\
IA3 & Infused Adapter by Inhibiting and Amplifying Inner Activations \\
ICT & Information and Communication Technology \\
JSONL & JSON Lines \\
LIME & Local Interpretable Model-Agnostic Explanations \\
LLM & Large Language Model \\
LoRA & Low-Rank Adaptation \\
NCSC & National Cyber Security Center \\
NF4 & NormalFloat 4-bit quantization \\
NLP & Natural Language Processing \\
OCR & Optical Character Recognition \\
OTP & One-Time Password \\
PEFT & Parameter-Efficient Fine-Tuning \\
QLoRA & Quantized Low-Rank Adaptation \\
RAM & Random Access Memory \\
SHA-256 & Secure Hash Algorithm 256-bit \\
SHAP & SHapley Additive exPlanations \\
SMS & Short Message Service \\
USTH & University of Science and Technology of Hanoi \\
URL & Uniform Resource Locator \\
VND & Vietnamese Dong \\
VRAM & Video Random Access Memory \\
XAI & Explainable Artificial Intelligence \\
\bottomrule
\end{longtable}

\clearpage
\chapter*{\listtablename}
\addcontentsline{toc}{chapter}{\listtablename}
{\singlespacing
\makeatletter\@starttoc{lot}\makeatother
}

\clearpage
\chapter*{\listfigurename}
\addcontentsline{toc}{chapter}{\listfigurename}
{\singlespacing
\makeatletter\@starttoc{lof}\makeatother
}

\clearpage
\chapter*{Abstract}
\addcontentsline{toc}{chapter}{Abstract}
This thesis presents a local-first pipeline for supervised four-class classification of Vietnamese financial phishing messages. For each message it returns a class label, a risk tier, the exact pieces of the message it used to decide, and safe next-step guidance, all on the user's own laptop, so a message containing bank details or a one-time password is never uploaded to a cloud service.

The first contribution is a 2{,}097-message dataset, built under strict rules, covering four classes: bank impersonation, Zalo social engineering, task scam, and benign. Several messages were written from each source article, so messages from the same article can look very similar to each other. All of them are kept in the same split, which stops the system from being tested on a message that is almost identical to one it already trained on. Every message is checked automatically for correct structure and correctly marked evidence spans, then scored by a separate model acting as a quality judge, and a Vietnamese-speaking reviewer checked 100 messages by hand, chosen evenly across all four classes. In a review of 324 messages whose labels looked doubtful, the reviewer agreed that 233 of them needed a different label. Only 57 of those corrections were used, because the remaining 176 messages all came from the same source article, and adding them would have made the dataset look more varied than it was. The same review found 240 Zalo messages written as third-person narration rather than as real chat messages, which were rebuilt offline from 60 preserved scenarios with no external API calls.

The second contribution is a training study on consumer hardware. Ordinary LoRA \cite{hu2022lora} did run, but on an 8~GiB laptop GPU it peaked at 7{,}902 of 8{,}151~MiB, left only 9~MiB free at its tightest, and projected roughly 18 hours of compute, making it impractical rather than impossible; no out-of-memory failure occurred. QLoRA was used instead, since its authors report that 4-bit quantized fine-tuning reaches 16-bit quality \cite{dettmers2023qlora}, making the switch a memory decision rather than an accuracy trade-off. Training completed locally for Qwen3-4B-Instruct-2507 with genuine NF4 QLoRA and for a non-quantized PhoBERT classification head, both on the same messages with one fixed seed. The selected Qwen model was exported as a Q8\_0 GGUF file. In one final test over 220 held-out messages, Qwen reached a macro F1 of 0.980493 and PhoBERT 0.990892, with no unreadable output from either. Each model was trained only once, so this gap describes what happened in this one test. It is not enough to prove that PhoBERT is always the better model. The messages themselves are generated from public warning articles rather than collected from real victims, so the prototype is not ready for real deployment until it has been evaluated on a real-world, human-labelled dataset.

\vspace{0.5em}
\noindent\textbf{Keywords:} Vietnamese phishing detection; local large language models; explainable artificial intelligence; QLoRA adaptation; GGUF inference; financial fraud prevention.

\cleardoublepage
\pagenumbering{arabic}
